% !PRINTED 409
\weberpage{409}
%<<CONT
\gesperrt{independent variables $x_0, x_1, \ldots, x_{n-1}$, one can
set}
\weq{4}{u = f(\vartheta),\ u_1 = f_1(\vartheta), \ldots, u_{\nu-1} = f_{\nu-1}(\vartheta)}
\gesperrt{where $f, f_1, \ldots, f_{\nu-1}$ are rational functions
of $\vartheta$ and}
\weq{5}{\vartheta = \chi(u, u_1, \ldots, u_{\nu-1}) = \psi(x_0, x_1, \ldots, x_{n-1})}
\gesperrt{is a rational function of the $u$, or also of the
$x$.}

For, by Theorem 1., if we replace
$x_0, x_1, \ldots, x_{n-1}$ by rational functions of a variable $t$,
in such a way that $z$ does not become independent of $t$, we can first represent the functions
$u_h$ by the formulae (4), (5). The relations
\[
u_h = f_h(\vartheta) = f_h[\chi(u, u_1, \ldots, u_{\nu-1})]
\]
are then identically satisfied with respect to $t$, and must therefore,
by Theorem 2., also be identical in the variables $x$,
q. e. d.

4.~\gesperrt{If of two variables $\vartheta, \vartheta_1$ each is a ratio\-
nal function of the other, then they are linear
functions of one another.}

For the proof we assume that
\weq{6}{\vartheta = \frac{\varphi(\vartheta_1)}{\psi(\vartheta_1)},\ \vartheta_1 = \frac{\varphi_1(\vartheta)}{\psi_1(\vartheta)},}
and understand by $\varphi, \psi$ two integral functions without common
 divisor, and likewise by $\varphi_1, \psi_1$. If we arrange the second of the
equations (6) in the form $\varphi_1(\vartheta) - \vartheta_1\psi_1(\vartheta) = 0$ according to powers
 of $\vartheta$, it may take the form:
\[
a_0\vartheta^m + a_1\vartheta^{m-1} + \cdots + a_{m-1}\vartheta + a_m = 0,
\]
in which the coefficients $a_0, a_1, \ldots, a_m$ are \gesperrt{integral linear functions}
 of $\vartheta_1$. If we substitute in it for $\vartheta$ the value
from the first of the equations (6), there follows the equation, identical
 with respect to $\vartheta_1$,
\[
a_0\varphi^m + a_1\varphi^{m-1}\psi + \cdots + a_{m-1}\varphi\psi^{m-1} + a_m\psi^m = 0.
\]
Accordingly $a_0\varphi^m$ must be divisible by $\psi$, and since $\varphi$
and $\psi$ are assumed to be relatively prime, $a_0$ must be divisible by $\psi$,
and hence $\psi$ constant or linear. In the same way we conclude
that $\varphi$ must be constant or linear, and since not both
%CONT>>
